% !TEX root = ../main.tex
% El recorrido del rastreador (Fase 1).
%
% Sustituye al diagrama de secuencia que habia antes, que ocupaba 0,9 de la
% altura de la pagina y aun asi habia que acercarse para leerlo. El tutor lo
% marco dos veces por eso mismo.
%
% El cambio no es solo de tamano: deja de ser una secuencia y pasa a ser el
% recorrido con sus bifurcaciones. Se gana legibilidad y se pierde detalle, y
% conviene tenerlo escrito: de las CUATRO bifurcaciones que dibujaba la
% version anterior aqui hay DOS. Se quedan fuera «guia con contenido frente a
% guia publicada vacia» y «grado simple frente a doble grado». No caben sin
% volver a la letra pequena que se queria evitar.
%
% El recuadro con borde grueso es la guia en PDF, que es donde se concentran
% las anomalias de la fuente: es la unica pieza que se sirve como binario y la
% unica con lista de secciones permitidas.
%
% Las cajas se encadenan con `below=of` y `left=of`: dar las alturas a mano ya
% provoco solapes en otras dos figuras, porque una caja crece con su texto y
% `minimum height` solo fija el minimo.
\begin{figure}[H]
  \centering
  % La nota va DEBAJO y no al lado a propósito. Puesta a la derecha estiraba el
  % dibujo hasta 19\,cm de ancho, y como `resizebox` conserva la proporción,
  % eso encogía toda la figura un 38\,\%: la letra volvía a quedarse pequeña,
  % que es justo lo que había que evitar.
  \resizebox{\textwidth}{!}{%
    \begin{tikzpicture}[
      paso/.style={nodoDiag, text width=5.0cm, minimum height=1.3cm},
      pasoFoco/.style={paso, draw=diagFocoBorde, fill=diagFocoFondo, very thick},
      rama/.style={nodoDiag, fill=white, densely dashed,
        text width=4.2cm, minimum height=1.3cm},
      decision/.style={diamond, draw=diagNodoBorde, fill=white, aspect=2.4,
        align=center, font=\small, inner sep=2pt},
      rotulo/.style={font=\scriptsize, inner sep=2pt},
      ]

      % ── La cadena de peticiones ────────────────────────────────────────────
      \node[paso] (n1)
      {Listado oficial de titulaciones\\{\scriptsize una sola petición}};

      \node[decision, below=0.7cm of n1] (d1)
      {¿Está en extinción?\\{\scriptsize se mira en el nombre}};

      \node[paso, below=0.7cm of d1] (n2)
      {Portada del grado\\{\scriptsize \emph{emite} \texttt{grado}}};

      \node[paso, below=0.7cm of n2] (n3)
      {Tabla de asignaturas\\{\scriptsize columnas por rótulo}\\{\scriptsize \emph{emite} \texttt{asignatura}}};

      \node[decision, below=0.7cm of n3] (d2)
      {¿Enlaza a su guía?\\{\scriptsize \texttt{href} real, nunca patrón}};

      \node[pasoFoco, below=0.7cm of d2] (n4)
      % El salto va a mano: sin él, «Descripción» parte en dos por la mitad.
      {Guía docente en PDF\\{\scriptsize solo «Resumen»}\\{\scriptsize y «Descripción de contenidos»}\\{\scriptsize \emph{emite} \texttt{guia}}};

      % ── Las ramas que no continúan ─────────────────────────────────────────
      \node[rama, left=1.4cm of d1] (n5)
      {Se descarta sin pedirla\\{\scriptsize no gasta una petición}};

      \node[rama, left=1.4cm of d2] (n6)
      {Queda anotada sin guía\\{\scriptsize \texttt{tiene\_guia: false}}};

      % ── La rama que sí continúa, y que sale de la portada ──────────────────
      \node[paso, right=1.4cm of n2] (n7)
      {Salidas profesionales\\{\scriptsize \emph{emite} \texttt{salidas}}};

      % ── Flechas ────────────────────────────────────────────────────────────
      % `(a) -- (b)` entre nodos recorta solo hasta sus bordes, así que no hace
      % falta calcular ni un punto.
      \draw[flechaDiag] (n1) -- (d1);
      \draw[flechaDiag] (d1) -- node[rotulo, above=2pt]{sí} (n5);
      \draw[flechaDiag] (d1) -- node[rotulo, right=2pt]{no} (n2);
      \draw[flechaDiag] (n2) -- (n7);
      \draw[flechaDiag] (n2) -- node[rotulo, right=2pt]{por cada grado} (n3);
      \draw[flechaDiag] (n3) -- node[rotulo, right=2pt]{por cada asignatura} (d2);
      \draw[flechaDiag] (d2) -- node[rotulo, above=2pt]{no} (n6);
      \draw[flechaDiag] (d2) -- node[rotulo, right=2pt]{sí} (n4);

      % ── Lo que el esquema por sí solo no dice ──────────────────────────────
      \node[rotulo, anchor=north, text width=12cm, align=center]
      at ([yshift=-0.8cm]n4.south)
      {Cada flecha hacia abajo es otra petición. Doce titulaciones,
        528 asignaturas y 288 guías: el recorrido se repite entero.};
    \end{tikzpicture}%
  }
  \caption[El recorrido del rastreador]{\textcolor{borrador2}{El recorrido del
      rastreador. No hace una única petición, sino que encadena una por cada
      tipo de página, y en cada caja está escrito qué registro emite. Las dos
      bifurcaciones son las que evitan trabajo inútil: una titulación en
      extinción se descarta \emph{antes} de pedirla, y una asignatura sin
      enlace de guía queda anotada como tal en lugar de desaparecer. Fuente:
      Elaboración propia.}}
  \label{fig:secuencia_scraping}
\end{figure}
